\documentclass[../main.tex]{subfiles}

\begin{document}

\ifSubfilesClassLoaded{

    %\tableofcontents
    
    \setcounter{chapter}{10}
}{}

\chapter{ HW11 }
代靖涵 25120222201319
\begin{exercise}{}{}
定义 \(\operatorname{GL}(n+1,\mathbb{R})\) 在 \(\mathbb{RP}^n\) 上的作用,
\[
\tau(A):\mathbb{RP}^n\to \mathbb{RP}^n,\qquad [x]\mapsto [Ax].
\]

a). 求证: 上式定义了 \(\operatorname{GL}(n+1,\mathbb{R})\) 在 \(\mathbb{RP}^n\) 上光滑的, 传递的左作用;

b). 上式中令 \(A\in \operatorname{GL}(n+1,\mathbb{C}),\ x\in \mathbb{CP}^n\), 求证:
\(\tau:\operatorname{GL}(n+1,\mathbb{C})\to \mathbb{CP}^n\) 是光滑的, 传递的左作用;

c). 将 \(S^{2n+1}\) 视为 \(\mathbb{C}^{n+1}\) 中的单位球面, \(S^1\) 在 \(S^{2n+1}\) 上的作用由下式给出 (Hopf 作用)
\[
z\cdot (w^1,\dots,w^{n+1})=(zw^1,\dots,zw^{n+1}).
\]
证明: 该作用是光滑的, 其轨道是 \(\mathbb{C}^{n+1}\) 中互不相交, 并集为 \(S^{2n+1}\) 的单位圆.


\end{exercise}

\begin{proof}
    \begin{enumerate}
        \item[a).]   对于任意的 \(  l \in \mathbb{R}   \), \(  x \in \mathbb{R} ^{n}  \), 我们有 \[
        [A\left( kx \right) ]= [A\left( kx \right) ]= [k \left( Ax \right) ]= [Ax]
        \] 因此 \(  \tau \left( A \right)   \)是良定义的映射.  考虑 \(  \mathbb{RP}^{n}  \)的图册 \[
        U_{k}= \left\{ [x_0: x_1:\cdots :x_{n}] : x_{k}\neq 0 \right\}
        \] 以及坐标映射\(   \varphi _{k}: U_{k}\to  \mathbb{R} ^{n} \)  \[
         \varphi _{k}\left( [x_0:x_1:\cdots :x_{n}] \right)=\left( \frac{x_0 }{x_{k} },\cdots , \frac{\hat{x_{k}} }{x_{k} } ,\cdots , \frac{x_{n} }{x_{k} }   \right)    
        \]存在  则对于 \(  A= \left( a_{ij} \right)\in  \operatorname{GL} \left( n+ 1 ,\mathbb{R} \right)    \), 对于 \(  x= [x_0:x_1,\cdots :x_{n}]  \)  , \(  \tau   \) 的坐标表示为 \[
       \begin{aligned}
  \left(  \left( a_{ij} \right),   \left( \frac{x_0 }{x_{k} },\cdots , \frac{\hat{x_{k}} }{x_{k} },\cdots , \frac{x_{n} }{x_{k} }    \right)  \right) & \mapsto  [\left( x_0,\cdots ,x_{n} \right) \left( a_{ij} \right) ]\\ 
        &= \left[ \sum _{i = 1}^{n+ 1}x_{i-1}a_{i 1} : \cdots :\sum _{i= 1}^{n+ 1}x_{i-1}a_{i \left( n+ 1 \right) } \right] \\ 
         &= \begin{pmatrix} 
             \frac{\sum _{i} ^{n+ 1}x_{i-1}a_{i1}}{\sum _{i = 1}^{n+ 1}x_{i-1}a_{i k} }  ,\cdots , \widehat{ \frac{\sum _{i= 1}^{n+ 1}x_{i-1}a_{i k} }{\sum _{i = 1}^{n+ 1}x_{i-1}a_{i k} } },\cdots , \frac{\sum _{i = 1}^{n+ 1}x_{i-1}a_{i \left( n+ 1 \right) } }{\sum _{i = 1}^{n+ 1}x_{i -1} a_{i k} } 
         \end{pmatrix} 
       \end{aligned}
        \] 对于任意点, 选取包含像点的坐标卡 $U_j$, 在局部邻域内, 该映射的坐标表示为有理分式且分母不为零, 此时上述表达式的每个分量都是关于 \(  x_0,\cdots ,\hat{x_{k}},\cdots ,x_{n}  \)和 \(  a_{ij},1\le i,j\le n+ 1  \) 的分式,  又 \(  \left( x_0,\cdots , \hat{x_{k}},\cdots ,x_{n} \right) \leftrightarrow \left( \frac{x_0 }{x_{k} },\cdots , \frac{\hat{x_{k}} }{x_{k} },\cdots ,\frac{x_{n} }{x_{k} }    \right)    \) 是 \(  U_{k}  \)上的微分同胚. 这就说明了 \(  \tau   \)是光滑的映射. 根据定义,  \[
        \begin{aligned}
        \tau \left( A  \right)  \left( \tau \left( B \right)\left( [x] \right)   \right)= \tau \left( A \right)[Bx]= [A \left( Bx \right) ]= [ABx]= \tau \left( Ab \right)[x]   
        \end{aligned}
        \] \[
        \tau \left( I_{n} \right)[x]= [x] 
        \]因此 \(  \tau   \)是 \(  \mathbb{RP}^{n}  \)上的一个左作用.  对于任意的 \(  [x]  , [y]\in \mathbb{RP}^{n}\). 由于 \(  \mathbb{R} ^{n}\to \mathbb{R} ^{n}, x\mapsto  Ax  \)是线性同构, 存在 \(  A_0\in  \operatorname{GL} \left( n+ 1,\mathbb{R}  \right)   \), 使得 \(  Ax= y  \), 从而  \(  \tau \left( A \right)[x]= [Ax]= [y]   \), 这表明 \(  \tau   \)的是传递的.
        \item [b).]   定义 \[
        \begin{aligned}
         F: \operatorname{GL} \left( n+ 1, \mathbb{C}  \right)\times \mathbb{C} ^{n+ 1}&\to \mathbb{C} ^{n+ 1} \\ 
          \left( A,x \right)&\mapsto  Ax
        \end{aligned}
        \]由于 \(  F  \)是 双线性的,  \(  F  \)是光滑映射. 为了计算 \(  \,\mathrm{d} F  \), 对于 \(  A\in \operatorname{GL} \left( n+ 1,\mathbb{C}  \right), x\in \mathbb{C} ^{n+ 1}   \), 取以 \(  \left( A, x \right)   \)为原点的曲线 \(   \gamma \left( t \right)= \left( A\left( t \right), x\left( t \right)   \right)    \), 使得\(   \gamma \left( 0 \right)= \left( A,x \right)    \).  则 \[
         \gamma ^{\prime} \left( 0 \right)= A^{\prime} \left( 0 \right)x\left( 0 \right)+ A\left( 0 \right)x^{\prime} \left( 0\right)= A^{\prime} \left( 0 \right)x+ Ax^{\prime} \left( 0 \right)      
        \]    因此 \[
        \,\mathrm{d} F_{\left( A,x \right) }\left( B,y \right)= Bx+ Ay 
        \] 特别地, \[
        \,\mathrm{d} F_{\left( A,x \right) }\left( 0,A^{-1} y \right)= y ,\forall y\in \mathbb{C} ^{n+ 1}
        \] 因此 \(  \,\mathrm{d} F_{\left( A,x \right) }  \)是满的. 这表明 \(  F  \)是淹没.  又注意到 \(  F  \)是满的, 因此 \(  F  \)是满的光滑淹没.  令 \(   \pi : \mathbb{C} ^{n+ 1}\to \mathbb{CP}^{n}  \)是商映射. 则 \(   \pi \circ F  \)也是满的光滑淹没.      对于任意的 \(  x\in \mathbb{C} ^{n+ 1}  \)和 \(   \lambda \in \mathbb{C}   \), \(   \pi \left( F\left(  \lambda x \right)  \right)=  \pi \left(  \lambda \left( Ax \right)  \right)=  \pi \left( Ax \right)=  \pi \left( F\left( x \right)  \right)      \)   .  \(  \operatorname{Id}_{\operatorname{GL} \left( n+ 1,\mathbb{C}  \right) } \times  \pi  \) 是商映射, 且 \(   \pi \circ F  \)在它的每个纤维 \(  \left\{ \left( A, \lambda x \right)  \right\}  \)上取常值 \(   \pi \left( F\left( x \right)  \right)   \). 因此 存在唯一的光滑映射 \(  \tilde{F}: \operatorname{GL} \left( n+ 1,\mathbb{C}  \right)\times \mathbb{CP}^{n}\to \mathbb{CP}^{n}   \)使得 \(  \tilde{F}\circ \left( \operatorname{Id}_{\operatorname{GL} \left( n+ 1 ,\mathbb{C} \right) }\times  \pi  \right)=  \pi \circ F   \). 由于 \(  \tau \circ \left( \operatorname{Id}_{\operatorname{GL} \left( n+ 1,\mathbb{C}  \right) }\times  \pi  \right)=  \pi \circ F   \), 故 \(  \tau  = \tilde{F}\)是光滑的.       故 \(  \tau   \)是光滑的作用. 与(a)类似地过程可知 \(  \tau   \)是一个传递的左作用. 
        \item [c).] 考虑映射 \(  F: \mathbb{C} ^{1}\times \mathbb{C} ^{n+ 1}\to \mathbb{C} ^{n+ 1}  \). \[
        F\left( z,\left( w^{1},\cdots ,w^{n+ 1} \right)  \right)= \left( zw^{1},\cdots ,zw^{n+ 1} \right)  
        \]  则 \(  F  \)是光滑映射. 令 \(  \iota_{n+ 1} : S^{2n+ 1}\hookrightarrow \mathbb{C} ^{n+ 1}  \)和 \(  \iota _{1}:S^{1}\hookrightarrow \mathbb{C} ^{1}  \) 是 典范嵌入.  则 \( \iota :=  \iota _{1}\times \iota _{n+ 1}: S^{1}\times S^{n+ 1}\hookrightarrow \mathbb{C} ^{1}\times \mathbb{C} ^{n+ 1}  \)也是嵌入. 因此 \(  F\circ \iota   \) 是光滑映射. 此外, 若 \(  z\in \mathbb{C} ^{1}, w\in \mathbb{C} ^{n+ 1}  \)使得 \(  \left| z \right|= 1   \) , \(  \left| w \right|= 1   \), 则 \[
        \left| \left( zw^{1},\cdots ,zw^{n+ 1} \right)  \right|= \left| z \right|\left| \left( w^{1},\cdots ,w^{n+ 1} \right)  \right|=\left| z \right|\left| w \right|=    1  
        \]   因此 \(  \left( F\circ \iota  \right)\left( S^{1}\times S^{2n+ 1} \right)\subseteq S^{2n+ 1}    \) 故 \(  F\circ \iota : S^{1}\times S^{2n+ 1}\to S^{2n+ 1}  \)是光滑映射. 它给出题述的Hopf作用, 因此 Hopf作用是光滑的作用. 任取 \(  w\in S^{2n+ 1}  \), \(  w  \)的轨道是 \(  \left\{ \left( zw^{1},\cdots ,z w^{n+ 1} \right): z \in \mathbb{C}   \right\}  \)    考虑映射 \[
        \left( zw^{1},\cdots ,zw^{n+ 1} \right)\to \left( z,0,\cdots ,0 \right)  
        \] 则 \[
        \begin{aligned}
      & \left<\left( z_1 w^{1},\cdots ,z_1w^{n+ 1} \right),\left( z_2w^{1},\cdots ,z_2 w^{n+ 1} \right)   \right> \\ 
       &= z_1w^{1} \overline{z_2w^{1}}+ \cdots +  z_1w^{n+ 1}\overline{z_2w^{n+ 1}}\\ 
        &= z_1 \bar{z_2} \left( \left| w^{1} \right|^{2}+ \cdots + \left| w^{n+ 1} \right|^{2}   \right)\\ 
         &= z_1  \overline{z}_{2} \\ 
          &= \left<\left( z_1,0,\cdots ,0 \right), \left( z_2,0,\cdots ,0 \right)   \right>
        \end{aligned}
        \]因此 \(  \left( zw^{1},\cdots ,z w^{n+ 1} \right)\to \left( z,0,\cdots ,0 \right)    \) 是等距同构. 同构像是 \(  S^{1}\times \left\{ 0 \right\}\times \cdots \times \left\{ 0 \right\}\simeq S^{1}  \). 因此Hopf作用的轨道是 单位圆. 对于任意的 \(  w_1,w_2\in  S^{2n+ 1}  \), \(  w_1\neq w_2  \), 都有 \(  \left( zw_1^{1},\cdots ,zw_1^{n+ 1} \right) \neq \left( zw_2^{1},\cdots ,zw_2^{n+ 1} \right) , \forall z\in \mathbb{C}   \).    因此轨道是互不相交的.  且对于任意的 \(  w\in S^{2n+ 1}  \), 取 \(  z= 1  \), 得到 \(  z\cdot w= w  \), 因此\(  S^{2n+ 1}  \)中的元素都含于某个轨道.    
    \end{enumerate}
    
\end{proof}


\begin{exercise}{}{}
特殊线性群 \(\operatorname{SL}(2,\mathbb{R})\) 通过 Möbius 变换作用在上半平面 \(\mathbb{H}=\{z:\operatorname{Im}(z)>0\}\) 上,
\[
\begin{pmatrix}
a & b\\
c & d
\end{pmatrix}\cdot z=\frac{az+b}{cz+d}.
\]

a). 求证: 该作用是传递的.

b). 求该作用下点 \(\sqrt{-1}\in \mathbb{H}\) 的稳定化子.


c). 让 \(G=   \operatorname{SL}\left( 2,\mathbb{R}  \right)   \)伴随地作用在自身上,  
\[
\operatorname{Ad}:G\to \operatorname{Aut}(G),\qquad g\mapsto \operatorname{Ad}_g,
\]
 \(B=\begin{pmatrix}0 & 1\\ -1 & 0\end{pmatrix}\) 的伴随轨道是
\[
O(B)=\{ABA^{-1}\mid A\in \operatorname{SL}(2,\mathbb{R})\}.
\]

 显式给出 \(\operatorname{SL}(2,\mathbb{R})\) 等变的微分同胚
\(f:\mathbb{H}\to O(B)\). 设 \(G\) 作用在光滑流形 \(M,N\) 上, 我们称 \(G\) 作用关于光滑映射 \(f:M\to N\) 是等变的 (equivariant), 如果满足
\[
f(g\cdot p)=g\cdot f(p)\quad \forall p\in M,g\in G.
\]
\end{exercise}

\begin{proof}
   \begin{enumerate}
    \item [a).]
    由于 \[
  \begin{pmatrix} 
        a&b\\ 
         c&d 
    \end{pmatrix} \cdot \left(   \begin{pmatrix} 
        e&f\\ 
         g&h 
    \end{pmatrix}\cdot z  \right)= \frac{a\frac{ez+ f }{gz+ h }+ b  }{c\frac{ez+ f }{gz+ h }+ d  }= \frac{\left( ae+ bg \right)z+ af+ bh  }{\left( ce+ gf \right)z+ \left( cf+ dh \right)   }= \left(\begin{pmatrix} 
        a&b\\ 
         c&d 
    \end{pmatrix} \begin{pmatrix} 
        e&f\\ 
         g&h 
    \end{pmatrix}   \right)\cdot z    
    \]因此该作用是一个左-李群作用.
    
    设 \(  z= x+ iy  \). 则 \[
   \begin{pmatrix} 
      \sqrt{y}&0\\ 
        0&\frac{1 }{\sqrt{y} }   
   \end{pmatrix} \cdot i=  iy
    \]  \[
    \begin{pmatrix} 
        1& x\\ 
         0& 1 
    \end{pmatrix} \cdot \left( iy \right)= x+ iy= z  
    \] 其中 \[
    \begin{pmatrix} 
        \sqrt{y}&0\\ 
         0&\frac{1 }{\sqrt{y} }  
    \end{pmatrix}, \begin{pmatrix} 
        1&x\\ 
         0&1 
    \end{pmatrix} \in \operatorname{SL}\left( 2,\mathbb{R}  \right)   
    \]因此对于任意的 \(  z \in \mathbb{H}  \), 存在 \(  A\in \operatorname{SL}\left( 2,\mathbb{R}  \right)   \), 使得 \(  A\cdot i= z  \), 则 \(  i= A^{-1} z  \)   . 对于任意的 \(  z_1,z_2\in \mathbb{H}  \), 设 \(  A_1\cdot i= z_1  \), \(  A_2\cdot i= z_2  \), 则 \[
    \left( A_2 A_1^{-1}  \right)z_1= z_2 
    \]   故该作用是传递的. 因此 \(  \mathbb{H}  \)是齐性的\(  \operatorname{SL}\left( 2,\mathbb{R}  \right)   \)-空间.  
    \item [b).] \[
    \mathrm{Stab}\left( i \right)= \left\{ A\in \operatorname{SL}\left( 2,\mathbb{R}  \right): Ai= i  \right\} 
    \]若 \(  A= \begin{pmatrix} 
        a&b\\ 
         c&d 
    \end{pmatrix}\in \mathrm{Stab}\left( i \right)    \), 则考虑到  \[
    \begin{pmatrix} 
        a&b\\ 
         c&d 
    \end{pmatrix}\cdot i= \frac{ai+ b }{ci+ d }  = i\iff  ai+ b= di-c\iff a= d, b= -c
    \] 且 \[
    \det \begin{pmatrix} 
        a&b\\ 
         -b&a 
    \end{pmatrix}= a^{2}+ b^{2} = 1
    \]因此 \(  A\in \mathrm{Stab}\left( i \right)   \), 当且仅当 \(  A  \)形如 \[
    A= \begin{pmatrix} 
        \cos  \theta & \sin  \theta \\ 
         -\sin  \theta &\cos  \theta  
    \end{pmatrix} 
    \]  因此 \[
    \mathrm{Stab}\left( i \right)= \operatorname{SO}\left( 2 \right) 
    \]
    \item [c).] 
    
    \[
    \mathrm{Stab}\left( B \right)= \left\{ A\in \operatorname{SL}\left( 2,\mathbb{R}  \right): ABA^{-1} = B  \right\} 
    \] 对于 \(  A= \begin{pmatrix} 
        a&b\\ 
         c&d 
    \end{pmatrix}   \) , \[
    AB= BA\iff \begin{pmatrix} 
        -b&a\\ 
         -d&c 
    \end{pmatrix}= \begin{pmatrix} 
        c&d\\ 
         -a&-b 
    \end{pmatrix}  \iff a= d, c= -b
    \]因此 \[
    \mathrm{Stab}\left( B \right)\simeq \mathrm{SO}\left( 2 \right)  
    \]定义 \[
    \begin{aligned}
    F: \operatorname{SL}\left( 2,\mathbb{R}  \right)\to  O\left( B \right) \\ 
     A\mapsto  ABA^{-1}    
    \end{aligned}
    \]则 \(  F  \)是一个光滑映射 , \(  \operatorname{Im}F= O\left( B \right)   \), 齐性空间的刻画定理 ,  \(  \operatorname{ker}F  = \mathrm{Stab}\left( B \right)\simeq \mathrm{SO}\left( 2 \right)  \)是一个闭的李子群, 进而给出\(  G  \)-等变的微分同胚   \[
   \tilde{F}: \operatorname{SL}\left( 2,\mathbb{R}  \right) /SO\left( 2 \right) \simeq   O\left( B \right)  ,\quad \tilde{F}\left( A \; \mathrm{SO}\left( 2 \right)  \right)= ABA^{-1}  
    \]
    另一方面, 由于 \(  \mathbb{H}  \)是齐性的 \(  \operatorname{SL}\left( 2,\mathbb{R}  \right)   \)-空间,   齐性空间的刻画定理给出一个 \(  \operatorname{SL}\left( 2,\mathbb{R}  \right)   \)-等变的微分同胚 \(   \varphi : G/\mathrm{Stab}\left( i \right)\simeq G/\mathrm{SO}\left( 2 \right)\to \mathbb{H}    \). \(   \varphi \left( A\; \mathrm{SO}\left( 2 \right)  \right)= A \cdot i   \)   .
    \[
    \operatorname{SL}\left( 2,\mathbb{R}  \right)/ \operatorname{SO}\left( 2 \right)\simeq   \mathbb{H}
    \]令 \(  f: \mathbb{H}\to O\left( B \right)   \), \(  f=  \tilde{F}\circ \varphi ^{-1}   \)   , \(  f\left( A\cdot i \right)= ABA^{-1}    \) . 具体地, 对于 \(  z= x+ iy  \) , 令  \[
    A= \begin{pmatrix} 
        1&x\\ 
         0&1 
    \end{pmatrix}\begin{pmatrix} 
        \sqrt{y}&0\\ 
         0&\frac{1 }{\sqrt{y} }  
    \end{pmatrix}  
    \] 则 (a)的过程表明 \(  A\cdot i= z  \). 因此 \(  f  \)的显式地表示为 \[
    \begin{aligned}
  \begin{aligned}
    f\left( x+ iy \right)&= ABA^{-1} =  \begin{pmatrix} 
        1&x\\ 
         0&1 
    \end{pmatrix}\begin{pmatrix} 
        \sqrt{y}&0\\ 
         0&\frac{1 }{\sqrt{y} }  
    \end{pmatrix}&\begin{pmatrix} 
        0&1\\ 
         -1&0 
    \end{pmatrix}\begin{pmatrix} 
        \frac{1 }{\sqrt{y} }&0\\ 
         0&\sqrt{y}  
    \end{pmatrix}    \begin{pmatrix} 
        1&-x\\ 
         0&1 
    \end{pmatrix} \\ 
     &= \begin{pmatrix} 
         1&x\\ 
          0&1 
     \end{pmatrix}\begin{pmatrix} 
           0&y\\ 
            -\frac{1 }{y }&0 
     \end{pmatrix}\begin{pmatrix} 
         1&-x\\ 
          0&1 
     \end{pmatrix}\\ 
      &= \begin{pmatrix} 
          -\frac{x }{y }& \frac{x^{2}+ y^{2} }{y }\\ 
           -\frac{1 }{y }&\frac{x }{y }     
      \end{pmatrix}      
  \end{aligned}  
    \end{aligned}
    \]   
   \end{enumerate}
   
\end{proof}
\begin{exercise}{}{}
设 $G=\operatorname{SL}(2,\mathbb{R})$. 令 $M=P_n$ 是变量为 $z_1,z_2$ 的 $n$ 次齐次多项式的线性空间.

a). 证明: 下式给出了 $\operatorname{SL}(2,\mathbb{R})$ 在 $M$ 上的一个光滑作用,
\[
\left(
\begin{pmatrix}
a & b\\
c & d
\end{pmatrix}
\cdot f
\right)(z_1,z_2):=f(az_1+cz_2,bz_1+dz_2).
\]

b). 证明:
\[
X=
\begin{pmatrix}
0 & 1\\
0 & 0
\end{pmatrix},\quad
Y=
\begin{pmatrix}
0 & 0\\
1 & 0
\end{pmatrix},\quad
H=
\begin{pmatrix}
1 & 0\\
0 & -1
\end{pmatrix}
\]
是 $\mathfrak{sl}(2,\mathbb{R})$ 的一组基, 且对应的无穷小作用 (infinitesimal action) 分别为
\[
X_M(f)=z_1\frac{\partial f}{\partial z_2},\quad
Y_M(f)=z_2\frac{\partial f}{\partial z_1},\quad
H_M(f)=z_1\frac{\partial f}{\partial z_1}-z_2\frac{\partial f}{\partial z_2}.
\]

一个群作用 $G\looparrowright M$ 的无穷小作用定义为
\[
\operatorname{d}r:\mathfrak{g}\to M
\]
\[
\operatorname{d}r(X)(p)=\left.\frac{\operatorname{d}}{\operatorname{d}t}\right|_{t=0}\exp(tX)\cdot p.
\]
\end{exercise}
\begin{proof}
    
\begin{enumerate}
    \item [a).] 考虑映射 \[
    \begin{aligned}
    F: \mathrm{M}\left( 2,\mathbb{R}  \right)\times \mathbb{R} [z_1,z_2]& \to \mathbb{R} [z_1,z_2] \\ 
    \left(  \begin{pmatrix} 
         a&b\\ 
          c&d 
     \end{pmatrix},f \right)\left( z_1,z_2 \right) &\mapsto   f\left( az_1+ cz_2,bz_1+ dz_2 \right)  
    \end{aligned}
    \]易见 \(  f\left( az_1+ cz_2,bz_1+ dz_2 \right)   \)每一项的系数都是关于 \(  a,b,c,d  \)  和 \(  f  \)的每一项的系数的多项式. 因此 \(  F  \)是光滑的映射. \(  P_{n}  \)是 \(  \mathbb{R} [z_1,z_2]  \)    的线性子空间, 从而是 \(  \mathbb{R} [z_1,z_2]  \)的嵌入子流形. 而 \(  \operatorname{SL}\left( 2,\mathbb{R}  \right)   \)是 \(  M\left( 2,\mathbb{R}  \right)   \)的嵌入子流形, 因此 \(  \operatorname{SL}\left( 2,\mathbb{R}  \right)\times P_{n}   \)是 \(  M\left( 2,\mathbb{R}  \right)   \times \mathbb{R} [z_1,z_2]\)     的嵌入子流形. 故 \(  F|_{\operatorname{SL}\left( 2,\mathbb{R}  \right)\times P_{n} }  \)也是光滑映射. 此外, 若 \(  f\in P_{n}  \), 设 \(  f\left( z_1,z_2 \right) = \sum _{k= 0}^{n}a_{k}z_1^{k} z_2^{n-k}  \), 则\[
      f\left( az_1+ cz_2, bz_1+ dz_2 \right)=    \sum _{k= 0}^{n}a_{k}\left( az_1+ cz_2 \right)^{k}\left( bz_1+ dz_2 \right)^{n-k}  
    \]也是 \(  n  \)次齐次的, 因此 \(  F\left( \operatorname{SL}\left( 2,\mathbb{R}  \right)\times P_{n}  \right)\subseteq P_{n}   \). 故 \(  F|_{\operatorname{SL}\left( 2,\mathbb{R}  \right)\times P_{n} }: \operatorname{SL}\left( 2,\mathbb{R}  \right)\times P_{n} \to P_{n}  \)是光滑映射. 此外,  \[
   \begin{aligned}
  &  \begin{pmatrix} 
        a_1&b_1\\ 
         c_1&d_1 
    \end{pmatrix}\cdot \left( \begin{pmatrix} 
       a_2&b_2\\ 
        c_2&d_2
    \end{pmatrix}\cdot f  \right)\left( z_1,z_2 \right)  \\ 
     & = \begin{pmatrix} 
        a_1&b_1\\ 
         c_1&d_1
    \end{pmatrix}\cdot f\left( a_2z_1+ c_2z_2,b_2z_1+ d_2z_2 \right)   \\ 
     &= f\left( \left( a_1a_2+ b_1c_2 \right)z_1+  \left(  c_1a_2+ d_1c_2 \right)z_2,\left( a_1b_2+ b_1d_2 \right)z_1+ \left( c_1b_2+ d_1d_2 \right)z_2   \right)\\ 
      &= \left( \begin{pmatrix} 
          a_1&b_1\\ 
           c_1&d_1 
      \end{pmatrix}\begin{pmatrix} 
          a_2&b_2\\ 
           c_2&d_2 
      \end{pmatrix}   \right)\cdot f\left( z_1,z_2 \right)   
   \end{aligned}
    \]   故光滑映射 \(  F|_{\operatorname{SL}\left( 2,\mathbb{R}  \right)\times P_{n} }: \operatorname{SL}\left( 2,\mathbb{R}  \right)\times P_{n} \to P_{n}  \) 给出了光滑的左作用.
    
    \item [b).] 易见 \(  X,Y  ,H\)是线性无关的, 且\(  \operatorname{dim}\mathfrak{sl}\left( 2,\mathbb{R}  \right)= \operatorname{dim}\operatorname{SL}\left( 2,\mathbb{R}  \right)= 3    \),  又 \[
    X,Y,H\in \mathfrak{sl}\left( 2,\mathbb{R}  \right)= \left\{ X \in M\left( 2,\mathbb{R}  \right): \operatorname{tr}\left( X \right)= 0   \right\} 
    \]  因此 \(  X,Y,H  \)是 \(  \mathfrak{sl}\left( 2,\mathbb{R}  \right)   \)的一组基.  \[
    X^{k}= O,\forall k> 1,\quad \exp \left( tX \right)= I+ tX= \begin{pmatrix} 
        1&t\\ 
         0&1 
    \end{pmatrix}  
    \]\[
    \left. \frac{\mathrm{d}}{\mathrm{d}t} \right|_{t= 0}\exp \left( tX \right)\cdot f \left( z_1,z_2 \right) =  \frac{\mathrm{d}}{\mathrm{d}t}|_{t= 0}f\left( z_1,tz_1+ z_2 \right)= z_1\frac{\partial f}{\partial z_2}\left( z_1,z_2 \right)  
    \]  因此 \[
    X_{M}\left( f \right)= z_1\frac{\partial f}{\partial z_2} 
    \] \[
    Y^{k}= O,\forall k> 1, \exp \left( tY \right)= I+ tY= \begin{pmatrix} 
        1&0\\ 
         t &1 
    \end{pmatrix}  
    \] \[
    \left. \frac{\mathrm{d}}{\mathrm{d}t} \right|_{t= 0}\exp \left( tY \right)\cdot f\left( z_1,z_2 \right)= \left. \frac{\mathrm{d}}{\mathrm{d}t} \right|_{t= 0}f\left( z_1+ tz_2,z_2 \right)= z_2\frac{\partial f}{\partial z_1}\left( z_1,z_2 \right)    
    \]因此 \[
    Y_{M}\left( f \right)= z_2\frac{\partial f}{\partial z_1} 
    \] \[
    \exp \left( tH \right)= \exp \left( \operatorname{diag} \left( t,-t \right)   \right)= \begin{pmatrix} 
        e^{t}&0\\ 
         0&e^{-t} 
    \end{pmatrix}   
    \] \[
    \left. \frac{\mathrm{d}}{\mathrm{d}t} \right|_{t= 0}\exp \left( tH \right)\cdot f\left( z_1,z_2 \right)= \left. \frac{\mathrm{d}}{\mathrm{d}t} \right|_{t= 0}f\left( e^{t}z_1, e^{-t}z_2 \right)= z_1\frac{\partial f }{\partial z_1}\left(z_1,z_2  \right)-z_2\frac{\partial f}{\partial z_2}\left( z_1,z_2 \right)     
    \]因此 \[
    H_{M}\left( f \right)= z_1\frac{\partial f}{\partial z_1}-z_2\frac{\partial f}{\partial z_2} 
    \]
\end{enumerate}

\end{proof}
\end{document}